\chapter{Resultados}

\noindent En este último capítulo presenta los resultados obtenidos en el desarrollo del proyecto. Cabe destacar que debido al acuerdo de confidencialidad existente entre AMAI y el CCD del ITAM sólo se comunicarán resultados que fueron publicados en referencia al Estudio XXII (2019-2020) del Estudio Anual de AMAI, a través del documento \emph{Estudio Anual de la Industria de Investigación de Mercados y Opinión Pública en México. Edición XXII (2019-2020).}\footnote{\url{https://www.amai.org/descargas/Estudio_Anual_AMAI_XXII_Edicion_2019_Comunicado_Medios_Junio_2020.pdf}}. El resto de resultados se ha compartido a manera de un reporte ejecutivo para las empresas participantes en el consabido estudio.

\section{Valor del mercado}

\begin{itemize}
    \item En 2019 la industria de investigación de mercados e inteligencia aplicada al consumidor en México alcanzó un valor anual de 7,664 millones de pesos.
    \item En esta edición hubo 71 empresas que participaron.
    \item Es posible hacer comparaciones directas de 61 empresas que participaron en la Edición XXII (2018-2019) y la Edición XXII (2019-2020) del Estudio Anual de AMAI.
    \item Las empresas comparables entre ambos periodos tuvieron un crecimiento orgánico de 0.81\%.
    \item Hablando de dichas empresas comparables, 56\% reportaron decrecimiento mientras que 44\% observaron un crecimiento.
\end{itemize}

\section{Distribución y volumen mercado}

\begin{itemize}
\item De acuerdo a los datos Edición XXII (2019-2020) del Estudio Anual de AMAI, en México sigue predominando la investigación de mercado basada en investigación de corte cuantitativo, que representa alrededor de 67\% de la facturación anual.
\item En 2019 se hicieron 9 mil proyectos de investigación e inteligencia de mercados, casi el mismo número que la edición anterior indicando que si bien no hubo un aumento tampoco hubo disminución en proyectos. 
\item Por lo que hace a unidades de investigación cualitativa se realizaron 9 mil sesiones de grupo, 13 mil entrevistas a profundidad y 4 mil observaciones etnográficas.
\item El número de entrevistas con informantes de proyectos de investigación en 2019 es de 7 millones de contactos efectivos. Las entrevistas personales -cara a cara más vía pública- continúan siendo las predominantes con 3.17 millones de contactos que representan aproximadamente el 43\% de la forma de contacto. En segundo lugar, se encuentran los puntos de contacto por Internet con 2.75 millones de contactos.
\end{itemize}


\section{Capital Humano}

\begin{itemize}
\item El capital humano para la Edición XXII (2019-2020) del Estudio Anual de AMAI estuvo conformado por casi 11 mil personas, que corresponde a una disminución de aproximadamente 61\% con respecto a la edición pasada\footnote{Aparentemente se debe al contexto de elecciones en el ejercicio previo, donde la mayor parte del personal activo se contrató de manera eventual}, e inclusive un poco más bajo que en entregas anteriores a la Edición XXI (2018-2019) donde había entre 12 y 15 mil personas en la industria.
\item Con respecto a la distribución del capital humano en términos de su género, es grato compartir que por primera vez, la proporción de hombres y mujeres es 1, es decir, por cada hombre hay una mujer. 

\item En esta edición los tipos de personal que en la edición anterior tuvieron más hombres que mujeres correspondientes a segmentos operativos y de producción, tuvieron una proporción de 1 y 1 respectivamente.
\end{itemize}

\section{Ranking de empresas}

A continuación se listan en esta tabla solamente aquellas empresas que expresamente manifestaron su aceptación de estar en este listado. En este sentido, este ranking no tiene el valor representativo de otros indicadores del Estudio Anual. 

Dicho listado se ha construido a partir de un orden, de mayor a menor, de facturación reportada; con lo cual la disposición de las empresas en orden no implica necesariamente una posición de mayor o menor calidad, prestigio o cualquier otro criterio que no sea el de la facturación reportada.

\begin{table}[]
\caption{Ranking de empresas participantes en Edición XXII (2019-2020) del Estudio Anual de AMAI (parte 1)}
\label{tab:ranking1}
\begin{tabular}{|l|l|}
\hline
Empresa                                                            & Rank \\ \hline
IPSOS                                                              & 1    \\ \hline
UPAX SA DE CV                                                      & 2    \\ \hline
DE LA RIVA                                                         & 3    \\ \hline
ESTADISTICA APLICADA                                               & 4    \\ \hline
SMART INDEX                                                        & 5    \\ \hline
O’ DE LA ROQUETTE                                                  & 6    \\ \hline
LEXIA INSIGHTS \& SOLUTIONS                                        & 7    \\ \hline
MARKETING GROUP                                                    & 8    \\ \hline
G DE VILLA Y ASOCIADOS SAPI DE CV                                  & 9    \\ \hline
GFK MEXICO                                                         & 10   \\ \hline
BRAIN                                                              & 11   \\ \hline
AAGA MARKETING                                                     & 12   \\ \hline
PEARSON                                                            & 13   \\ \hline
NETQUEST                                                           & 14   \\ \hline
PHENOMA                                                            & 15   \\ \hline
SERTA MARKETING INTELLIGENCE PATTNER                               & 16   \\ \hline
EL INSTITUTO DE INVESTIGACIONES SOCIALES SC                        & 17   \\ \hline
\parbox[t]{11cm}{NODO INVESTIGACION ESTRATEGICA SUASOR CONSULTORES}                  & 18   \\ \hline
SUASOR CONSULTORES                                                 & 19   \\ \hline
\parbox[t]{11cm}{ESTUDIOS DE VARIABLES DEL MERCADO SC BERUMEN Y ASOCIADOS S A DE CV} & 20   \\ \hline
\end{tabular}
\end{table}


\begin{table}[]
\caption{Ranking de empresas participantes en Edición XXII (2019-2020) del Estudio Anual de AMAI (parte 2)}
\label{tab:ranking2}
\begin{tabular}{|l|l|}
\hline
Empresa                                                      & Rank \\ \hline
BERUMEN Y ASOCIADOS S A DE CV                                & 21   \\ \hline
PQR PLANNING QUANT                                           & 22   \\ \hline
FOCUS INVESTIGACION DE MERCADOS                              & 23   \\ \hline
INTEGRACION TOTAL                                            & 24   \\ \hline
MARES CONSUMER INTELLIGENCE                                  & 25   \\ \hline
INMEGA DESARROLLO DE NEGOCIOS                                & 26   \\ \hline
MERCAEI                                                      & 27   \\ \hline
QUESTIONPRO LATINOAMERICA                                    & 28   \\ \hline
\parbox[t]{11cm}{CUARTEL GENERAL DE INVESTIGACION DE MERCADOS MASTER RESEARCH} & 29   \\ \hline
MASTER RESEARCH                                              & 30   \\ \hline
PSYMA LATINA                                                 & 31   \\ \hline
FACTA RESEARCH                                               & 32   \\ \hline
EVIDENS CONSULTORIA DE MARCA                                 & 33   \\ \hline
GAIN DYNAMICS RESEARCH                                       & 34   \\ \hline
ENKOLL                                                       & 35   \\ \hline
DINAMIA                                                      & 36   \\ \hline
CINCO                                                        & 37   \\ \hline
INMERSA MARKETING GROUP                                      & 38   \\ \hline
ACSI RESEARCH                                                & 39   \\ \hline
IMAAC MARKETING GROUP                                        & 40   \\ \hline
\end{tabular}
\end{table}

\begin{center}
\begin{table}[h]
\caption{Ranking de empresas participantes en Edición XXII (2019-2020) del Estudio Anual de AMAI (parte 3)}
\label{tab:ranking3}
\begin{tabular}{|l|l|}
\hline
Empresa                   & Rank \\ \hline
SINCRONIA                 & 41   \\ \hline
ISOPOMER                  & 42   \\ \hline
MEBA                      & 43   \\ \hline
WISUM                     & 44   \\ \hline
GAUSSC                    & 45   \\ \hline
NEUROMARKETING            & 46   \\ \hline
SEGMENTOS RESEARCH        & 47   \\ \hline
SURVEY                    & 48   \\ \hline
OVALBOX                   & 49   \\ \hline
SEMIOSFERA CONSULTING S C \hspace{4cm} & 50   \\ \hline
\end{tabular}
\end{table}
\end{center}

\newpage
% Gráfica con dos colores a mano

\section{Documentos generados para AMAI}

Por otra parte, cabe destacar que como resultado del proyecto recién descrito, se entregaron dicha asociación una serie de documentos que contienen los resultados generados con el producto de datos diseñado para la Edicion XXII del Estudio Anual de AMAI:

\begin{enumerate}
    \item \emph{Estudio Anual de la Industria de Investigación de Mercados y Opinión Pública en México. Edición XXII (2019-2020). Nota metodológica.}\footnote{Disponible a través de la dirección electrónica: \url{http://amai.org/descargas/reporteMetodologicoCliente270520.pdf}}
    \item \emph{Estudio Anual de la Industria de Investigación de Mercados y Opinión Pública en México. Edición (2019-2020). Comunicado a Medios.}\footnote{Véase \url{https://www.amai.org/descargas/Estudio_Anual_AMAI_XXII_Edicion_2019_Comunicado_Medios_Junio_2020.pdf}}
    \item \emph{Estudio Anual de la Industria de Investigación de Mercados y Opinión Pública en México. Edición (2019-2020). Reporte a Participantes}
\end{enumerate}

Sobre éste último documento, cabe destacar que debido al acuerdo de confidencialidad existente entre AMAI y el CCD los resultados se reservaron para participantes en el estudio debido al carácter confidencial de la información, sin embargo su contenido versa sobre los elementos de valor obtenidos en el estudio con referencia a los puntos descritos en la sección \ref{obj:bussiness}.